\begin{tikzpicture}[
  >={Latex[length=2.5mm,width=1.8mm]},
  every node/.style={font=\small},
  bvec/.style={->,line width=1pt,bblue}
]

% --- Sinistra: B entra ed esce ---
\begin{scope}
  \fill[blobpink!40,draw=blobpink!70!black]
    plot [smooth cycle,tension=0.85] coordinates {
      (-1.8,0.3) (-1.0,2.0) (0.5,2.2) (1.6,1.5) (2.0,0.2)
      (1.5,-1.2) (0.0,-1.5) (-1.2,-1.2) (-2.0,-0.3)
    };

  % Cilindretto gaussiano Σ
  \fill[surfblue,opacity=0.5,draw=bblue] (-0.4,-0.6) rectangle (0.4,0.6);
  \draw[bblue] (0,0.6) ellipse (0.4 and 0.1);
  \draw[bblue,dashed] (0,-0.6) ellipse (0.4 and 0.1);

  % Linee di B (curve blu verticali)
  \foreach \x in {-1.2,-0.6,0.0,0.6,1.2}{
    \draw[bvec,smooth] (\x,-1.8) .. controls (\x,-0.5) and (\x,0.5) .. (\x,1.0);
    \draw[bvec] (\x,1.0) -- (\x,2.5);
  }

  % Piccoli loop interni (correnti amperiane)
  \foreach \y in {-0.2,0.3,0.8}{
    \draw[bblue,thin] (-0.3,\y) ellipse (0.2 and 0.08);
  }

  \node[bblue,above] at (0.6,2.5) {$\mathbf{B}$};
\end{scope}

% --- Destra: B esce da corpo diverso ---
\begin{scope}[xshift=5.5cm]
  \fill[blobpink!40,draw=blobpink!70!black]
    plot [smooth cycle,tension=0.85] coordinates {
      (-1.5,0.5) (-0.6,1.8) (0.8,2.0) (1.8,1.3) (2.2,0.0)
      (1.5,-1.2) (0.0,-1.5) (-1.0,-1.0) (-1.8,-0.2)
    };

  \fill[surfblue,opacity=0.5,draw=bblue] (-0.3,-0.5) rectangle (0.3,0.5);
  \draw[bblue] (0,0.5) ellipse (0.3 and 0.08);
  \draw[bblue,dashed] (0,-0.5) ellipse (0.3 and 0.08);

  \foreach \x in {-0.8,-0.2,0.4,1.0}{
    \draw[bvec,smooth] (\x,-1.8) .. controls (\x,-0.3) and (\x,0.3) .. (\x,0.8);
    \draw[bvec] (\x,0.8) -- (\x,2.5);
  }

  \node[bblue,right] at (1.5,2.2) {$\mathbf{B}$};
  \node[right] at (2.0,0.6) {$\Sigma$};
\end{scope}

\end{tikzpicture}
